\documentclass[11pt,a4paper]{article}

% ---- Packages ----
\usepackage[utf8]{inputenc}
\usepackage[T1]{fontenc}
\usepackage[english]{babel}
\usepackage{geometry}
\geometry{margin=2.5cm}
\usepackage{booktabs}
\usepackage{enumitem}
\usepackage{amsmath,amssymb}
\usepackage{hyperref}
\usepackage{xcolor}
\usepackage{fancyhdr}
\usepackage{titlesec}
\usepackage{listings}
\usepackage{tabularx}
\usepackage{float}

% ---- Colors ----
\definecolor{pts}{RGB}{180,60,60}

% ---- Code listings ----
\lstset{
  language=Python,
  basicstyle=\ttfamily\small,
  keywordstyle=\bfseries,
  commentstyle=\color{gray}\itshape,
  stringstyle=\color{green!40!black},
  breaklines=true,
  frame=single,
  rulecolor=\color{gray!50},
  backgroundcolor=\color{gray!5},
  numbers=left,
  numberstyle=\tiny\color{gray},
  tabsize=4,
}

% ---- Header/Footer ----
\pagestyle{fancy}
\fancyhf{}
\fancyhead[L]{\small INFO8003 -- Reinforcement Learning}
\fancyhead[R]{\small Project 2025--2026}
\fancyfoot[C]{\thepage}
\renewcommand{\headrulewidth}{0.4pt}

% ---- Section formatting ----
\titleformat{\section}{\Large\bfseries}{\thesection}{1em}{}
\titleformat{\subsection}{\large\bfseries}{\thesubsection}{1em}{}

% ---- Hyperlinks ----
\hypersetup{
  colorlinks=true,
  linkcolor=black,
  urlcolor=blue!60!black,
  citecolor=black,
}

% ---- Points macro ----
\newcommand{\pts}[1]{\hfill{\normalfont\normalsize\color{pts} [#1/20]}}

\begin{document}

% ====================================================================
% TITLE PAGE
% ====================================================================
\begin{titlepage}
\centering
\vspace*{3cm}

{\Huge\bfseries Reinforcement Learning Project}\\[0.8cm]
{\LARGE UAV Control via Deep Reinforcement Learning}\\[0.5cm]
{\Large PyFlyt Simulator}\\[2cm]

\begin{tabular}{ll}
\toprule
\textbf{Course}     & INFO8003 -- Reinforcement Learning \\
\textbf{Groups}     & 2 students \\
\textbf{Simulator}  & PyFlyt v0.29+ \cite{tai2023pyflyt} \\
\bottomrule
\end{tabular}

\vspace{2cm}
\textit{Train, evaluate, and compare deep RL agents for drone control.\\
Compete in a dogfight tournament.}

\vfill
{\large Academic Year 2025--2026}
\end{titlepage}

\tableofcontents
\newpage

% ====================================================================
% LLM POLICY
% ====================================================================
\section*{Use of LLMs}

You may use large language models (ChatGPT, Claude, Copilot, etc.) to help write code. However:
\begin{itemize}[nosep]
  \item You \textbf{must understand every line of code} you submit.
  \item \textbf{Questions will be asked} at the oral presentation and the exam. You must be able to explain your implementation, your training pipeline, and the algorithms you used.
  \item If you cannot explain your code, it will be considered as not done.
\end{itemize}

% ====================================================================
% 1. ENVIRONMENTS & MDP FORMALIZATION
% ====================================================================
\section{Environments \& MDP Formalization\pts{3}}
\label{sec:environments}

All environments use PyFlyt \cite{tai2023pyflyt}, a Gymnasium-compatible UAV simulator built on PyBullet. They follow the standard Gymnasium API:

\begin{lstlisting}[language=Python, numbers=none]
import gymnasium
env = gymnasium.make("PyFlyt/QuadX-Hover-v4", flight_mode=0)
obs, info = env.reset(seed=42)        # obs: np.ndarray
action = env.action_space.sample()    # action: np.ndarray(4,) in [-1, 1]
obs, reward, terminated, truncated, info = env.step(action)
env.close()
\end{lstlisting}

Any RL library that works with Gymnasium can be used. The 4D continuous action space is shared across all environments; its meaning depends on the \texttt{flight\_mode}.

In your report, \textbf{formalize each environment as an MDP}. Explore the environments in code to understand what the observations and rewards look like.

\subsection{Flight Modes}

\begin{table}[H]
\centering
\begin{tabularx}{\textwidth}{clXl}
\toprule
\textbf{Mode} & \textbf{Action semantics} & \textbf{Description} & \textbf{Difficulty} \\
\midrule
$-1$ & $(m_1, m_2, m_3, m_4)$ & Direct motor PWM & Hardest \\
$0$  & $(\dot{p}, \dot{q}, \dot{r}, T)$ & Angular velocity + thrust & Hard \\
$4$  & $(u, v, \dot{r}, z)$ & Body-frame velocity + yaw rate + height & Medium \\
$6$  & $(\dot{x}, \dot{y}, \dot{r}, \dot{z})$ & Ground velocity + yaw rate + vertical vel. & Easy \\
$7$  & $(x, y, r, z)$ & Position + yaw + height (full PID) & Easiest \\
\bottomrule
\end{tabularx}
\caption{Selected flight modes. Higher modes abstract away low-level control via cascaded PID controllers.}
\label{tab:flight_modes}
\end{table}

\subsection{QuadX-Hover-v4 (Stabilization)}

Maintain a quadrotor at a fixed target position with stable orientation.

\subsection{QuadX-Waypoints-v4 (Navigation)}

Navigate through randomly placed 3D waypoints. The default PyFlyt parameters are very restrictive. The provided code overrides them in \texttt{env\_config.py}. Read and understand those overrides.

\begin{lstlisting}[language=Python, numbers=none]
env = gymnasium.make("PyFlyt/QuadX-Waypoints-v4", flight_mode=6,
    goal_reach_distance=4.0, flight_dome_size=150.0,
    max_duration_seconds=120.0, num_targets=4)
\end{lstlisting}

This environment returns a \texttt{Dict} observation space. Most RL libraries need a flat vector. The provided \texttt{FlattenWaypointEnv} wrapper handles this.

\subsection{Dogfight -- MAFixedwingDogfightEnvV2}

1v1 fixed-wing aerial combat. This is a multi-agent environment (PettingZoo). The provided \texttt{DogfightSelfPlayEnv} wrapper turns it into a single-agent Gymnasium environment for training.

% ====================================================================
% 2. ALGORITHMS
% ====================================================================
\section{Algorithms}
\label{sec:algorithms}

You may use \textbf{any RL algorithm} and \textbf{any library}. You must train and compare \textbf{at least two algorithms}.

In your report, describe each algorithm: its main idea, key properties (on-policy vs off-policy, exploration mechanism, etc.), and why you chose it. Explain things in your own words.

Your trained model must expose the following interface (used by the evaluation scripts and the tournament):

\begin{lstlisting}[language=Python, numbers=none]
action, info = model.predict(obs, deterministic=True)
# obs:    np.ndarray  -- observation from the environment
# action: np.ndarray  -- action to take (shape depends on env)
# info:   any         -- optional extra info (can be {})
\end{lstlisting}

% ====================================================================
% 3. EXPERIMENTAL PROTOCOL
% ====================================================================
\section{Experiments \& Results\pts{4}}
\label{sec:protocol}

I encourage to log your runs to \href{https://wandb.ai}{Weights \& Biases}: episode reward, evaluation reward (deterministic policy), and policy videos.

Use appropriate statistical tools to compare your algorithms. Do not draw big conclusions on a single run of training. The quality of your evaluation matters.

\textit{Hint}: Agarwal et al.\ \cite{agarwal2021rliable} discuss best practices for statistical evaluation in deep RL. Worth reading.

% ====================================================================
% 4. ANALYSIS & DISCUSSION
% ====================================================================
\section{Analysis \& Discussion\pts{5}}
\label{sec:analysis}

This section carries the most weight. We care more about your understanding than about high scores. Your report should cover:

\begin{itemize}[nosep]
  \item Why does one algorithm work better than the other on a given environment?
  \item Effect of flight mode on waypoint performance: why is mode 0 harder than mode 6?
  \item Failure analysis: what causes crashes, reward plateaus, or policy collapse?
  \item What would you do differently with more compute?
\end{itemize}

% ====================================================================
% 5. IMPLEMENTATION
% ====================================================================
\section{Implementation\pts{4}}
\label{sec:implementation}

\begin{itemize}[nosep]
  \item Clean, readable code with \texttt{requirements.txt} and \texttt{README.md}.
  \item Proper use of Weights \& Biases for experiment tracking.
  \item Reproducible results: seeded experiments, documented hyperparameters.
  \item Any RL library is fine, implementing your own algorithm or exploring an algorithm proposed in papers is encouraged and rewarded.
\end{itemize}

% ====================================================================
% 6. DOGFIGHT TOURNAMENT
% ====================================================================
\section{Dogfight Tournament\pts{2}}
\label{sec:tournament}

\subsection{Rules}

Each group may submit \textbf{as many relevant agents as they want}. Test them before submitting. Motivate how you want to obtain the best agent and what you have done.

The tournament environment is fixed:
\begin{lstlisting}[language=Python, numbers=none]
MAFixedwingDogfightEnvV2(
    team_size=1, assisted_flight=True,
    flatten_observation=True, max_duration_seconds=60, agent_hz=30,
)
\end{lstlisting}

\subsection{Submission Format}

Submit Python files named \texttt{groupXX\_name.py} (e.g., \texttt{group03\_aggressive.py}) with a \texttt{load\_model()} function. The only requirement is the \texttt{predict()} interface. See \texttt{scripts/submission\_template.py}.

\subsection{Scoring}

Agents play a round-robin tournament, Elo-rated. Ranking bonus: top-3 (+1pt), top-5 (+0.5pt), based on your best agent only.

\subsection{Leaderboard}

A live leaderboard will be made available in the coming weeks. You will be able to submit agents and track your ranking before the final deadline.

% ====================================================================
% 7. BONUS / INNOVATION
% ====================================================================
\section{Bonus / Innovation\pts{2}}
\label{sec:bonus}

Extra credit for going beyond the required experiments: hyperparameter ablation, reward shaping, curriculum learning, architecture exploration, additional algorithms, etc.

% ====================================================================
% 8. DELIVERABLES
% ====================================================================
\section{Deliverables}
\label{sec:deliverables}

You are expected to submit a complete zipped file with the following on the Montefiore submission platform before the \textbf{12th May}.
\begin{enumerate}[nosep]
  \item \textbf{Code repository} (Git): training scripts, wrappers, analysis. Include \texttt{requirements.txt} and \texttt{README.md}.
  \item \textbf{Report} (PDF): MDP formalization, algorithm description, experimental setup, results with statistical analysis, flight mode comparison, dogfight strategy, discussion, conclusion.
  \item \textbf{Wandb project link (optional but recommended)}: public or shared.
  \item \textbf{Trained models}: best checkpoints per (algorithm, environment) pair.
  \item \textbf{Tournament submissions}: \texttt{groupXX\_name.py} files.
\end{enumerate}

% ====================================================================
% 9. PROVIDED CODE
% ====================================================================
\section{Provided Code}
\label{sec:starter}

The repository contains environment configuration and evaluation scripts.

\begin{table}[H]
\centering
\begin{tabularx}{\textwidth}{lX}
\toprule
\textbf{Script} & \textbf{Purpose} \\
\midrule
\texttt{scripts/env\_config.py} & Environment parameters (waypoint overrides) \\
\texttt{scripts/wrappers.py} & \texttt{FlattenWaypointEnv}: flattens dict observations to a fixed-size vector \\
\texttt{scripts/dogfight\_wrapper.py} & \texttt{DogfightSelfPlayEnv}: wraps multi-agent dogfight as single-agent Gymnasium env \\
\texttt{scripts/evaluate.py} & Evaluate any model (\texttt{.py} module or \texttt{.zip} checkpoint) \\
\texttt{scripts/tournament.py} & Dogfight tournament (used for grading) \\
\texttt{scripts/submission\_template.py} & Tournament submission template \\
\bottomrule
\end{tabularx}
\end{table}

\texttt{evaluate.py} and \texttt{tournament.py} are what we use to grade your work. Make sure your models run correctly with them before submitting.

% ====================================================================
% REFERENCES
% ====================================================================
\newpage
\bibliographystyle{plain}
\begin{thebibliography}{9}

\bibitem{tai2023pyflyt}
Tai, J.~J. et al. (2023).
\textit{PyFlyt -- UAV Simulation Environments for Reinforcement Learning Research}.
arXiv:2304.01305.

\bibitem{agarwal2021rliable}
Agarwal, R. et al. (2021).
\textit{Deep Reinforcement Learning at the Edge of the Statistical Precipice}.
NeurIPS 2021.

\end{thebibliography}

\end{document}
